\documentclass{beamer}

\hypersetup{unicode}

\usepackage[utf8]{inputenc}
\usepackage[T1]{fontenc}
\usepackage{textcomp}
\usepackage{lmodern}
\usepackage[croatian]{babel}
\usepackage{csquotes}
\MakeOuterQuote{"}

\usepackage{thmtools}

\declaretheorem{teorem}
\declaretheorem[sibling=teorem, style=plain]{lema}
\declaretheorem[style=remark, sibling=teorem]{napomena}
\declaretheorem[style=remark, sibling=teorem]{komentar}
\declaretheorem[style=definition, sibling=teorem]{definicija}
\declaretheorem[sibling=teorem]{korolar}
\declaretheorem[style=remark, sibling=teorem]{zadatak}
\declaretheorem[sibling=teorem]{propozicija}
\declaretheorem[style = example]{primjer}

\usepackage[style=numeric, sorting=none, giveninits]{biblatex}
\addbibresource{literatura-prez.bib}

\usepackage{tikz}
\usetikzlibrary{arrows}
\newcommand{\midarrow}{\tikz \draw[-triangle 90] (0,0) -- +(.1,0);}

\definecolor{nixblue}{HTML}{5277C3}
\definecolor{nixviolet}{HTML}{7C6BC4}
\definecolor{nixlight}{HTML}{7EBAE4}
\definecolor{white}{HTML}{ffffff}
\definecolor{black}{HTML}{000000}
\definecolor{lblack}{HTML}{222222}

\usetheme{Madrid}
\usecolortheme{dolphin}
\setbeamercolor{structure}{fg=nixblue}
\setbeamercolor{frametitle}{bg=nixblue, fg=white}
\setbeamercolor{palette primary}{bg=nixblue, fg=white}
\setbeamercolor{palette secondary}{bg=nixblue, fg=white}
\setbeamercolor{palette tertiary}{bg=nixblue, fg=white}
\setbeamercolor{palette quaternary}{bg=nixblue, fg=white}
\setbeamertemplate{navigation symbols}{}

\usepackage{listings}
\lstdefinestyle{mystyle}{
	backgroundcolor=\color{white},
	commentstyle=\color{nixviolet},
	keywordstyle=\color{nixblue},
	stringstyle=\color{nixviolet},
	basicstyle=\ttfamily\footnotesize,
	breakatwhitespace=false,
	breaklines=true,
	captionpos=b,
	keepspaces=true,
	numbers=none,
	showspaces=false,
	showstringspaces=false,
	showtabs=false,
	tabsize=2
}
\lstdefinelanguage{nix}{
	keywords={let, in, with, rec, inherit, import, if, then, else, assert, true, false, null},
	sensitive=true,
	comment=[l]{\#},
	morecomment=[s]{/*}{*/}
}
\lstset{
	style=mystyle,
	extendedchars=true,
	literate={→}{{\textrightarrow}}1 {←}{{\textleftarrow}}1
	         {↓}{{\textdownarrow}}1 {↑}{{\textuparrow}}1
	         {ć}{{\'c}}1 {č}{{\v c}}1 {š}{{\v s}}1 {ž}{{\v z}}1 {đ}{{\dj}}1
	         {Ć}{{\'C}}1 {Č}{{\v C}}1 {Š}{{\v S}}1 {Ž}{{\v Z}}1 {Đ}{{\DJ}}1
}

\usepackage{booktabs}
\usepackage{array}

\renewcommand{\mathrm}[1]{\mathsf{#1}}
\newcommand{\ub}{\rightarrow \infty}
\newcommand{\un}{\rightarrow 0}
\newcommand{\norm}[1]{\left\lVert#1\right\rVert}
\newcommand{\dotnorm}{\left\lVert \, \cdot \, \right\rVert}
\newcommand{\given}{\, \vert \,}
\newcommand{\st}{\, \colon \,}
\newcommand{\D}{\,\mathrm d}
\newcommand{\ani}{{^\perp}}
\newcommand{\ran}{\mathrm{ran}\,}
\newcommand{\bop}{\mathrm{B}}
\newcommand{\skp}[2]{\left\langle {#1}, {#2}  \right\rangle}
\newcommand{\jpod}{\stackrel{\mathrm d}{=}}
\newcommand{\jgs}{\stackrel{\mathrm{g.s.}}{=\joinrel=}}
\newcommand{\stackrell}[2]{\stackrel{\scriptstyle {#1}}{#2}}
\newcommand{\konvp}{\stackrel{\mathbb P}{\longrightarrow}}
\newcommand{\konvw}{\stackrel{\mathrm w}{\longrightarrow}}
\newcommand{\konvd}{\stackrel{\mathrm d}{\longrightarrow}}
\newcommand{\konvgs}{\stackrel{\mathrm{g.s.}}{\longrightarrow}}
\DeclareMathOperator{\card}{card}
\DeclareMathOperator{\var}{Var}
\DeclareMathOperator{\supp}{supp}
\DeclareMathOperator{\im}{im}
\DeclareMathOperator{\olspan}{\overline{span}}
\DeclareMathOperator{\re}{Re}
\DeclareMathOperator{\sgn}{sgn}

\newcommand{\ol}[1]{\overline{#1}}
\newcommand{\wh}[1]{\widehat{#1}}
\newcommand{\wt}[1]{\widetilde{#1}}
\renewcommand{\th}[1]{\hat{#1}}

\newcommand{\R}{\mathbb{R}}
\newcommand{\N}{\mathbb{N}}
\newcommand{\Z}{\mathbb{Z}}
\newcommand{\Q}{\mathbb{Q}}
\newcommand{\C}{\mathbb{C}}
\renewcommand{\P}{\mathbb{P}}
\newcommand{\Ptrans}{\P}
\newcommand{\E}{\mathbb{E}}
\def\B{\mathrm{B}}
\def\K{\mathrm{K}}
\def\M{\mathrm{M}}
\def\CC{\mathrm{C}}
\def\RR{\mathrm{R}}
\def\H{\mathcal{H}}
\def\K{\mathcal{K}}
\def\S{\mathcal{S}}
\def\T{\mathcal{T}}
\def\U{\mathcal{U}}

\newcommand{\cz}{Calder\' on--Zygmund}
\newcommand{\holder}{H\" older}
\newcommand{\cadlag}{c\` adl\` ag}
\newcommand{\levy}{L\' evy}
\newcommand{\ito}{It\= o}
\newcommand{\dnorm}{\norm{\, \cdot \,}}
\newcommand{\mathgs}{\ \ \mathrm{g.s.}}
\def\L{\mathrm{L}}
\def\w{\mathrm{w}}
\def\mrC{\mathrm{C}}
\def\mrc{\mathrm{c}}
\newcommand{\abs}[1]{\left| {#1} \right|}
\newcommand{\advisor}{Voditelj rada: prof.\ dr.\ sc.\ Nikola Sandrić}

\newcommand{\pitanje}[1]{%
	\begin{center}
		\vspace{1.2em}
		{\Large\itshape\usebeamercolor[fg]{structure}#1}
	\end{center}
}

\newcommand{\demo}[1]{%
	\begin{center}
		\vspace{0.8em}
		{\small\textcolor{nixblue}{\textsc{demo}}\quad\itshape\textcolor{black!70}{#1}}
	\end{center}
}

\let\etoolboxforlistloop\forlistloop
\usepackage{autonum}
\let\forlistloop\etoolboxforlistloop

\usepackage{enumitem}
\setitemize{label=\usebeamerfont*{itemize item}%
    \usebeamercolor[fg]{itemize item}
    \usebeamertemplate{itemize item}}
\makeatletter
\providecommand*{\blx@noerroretextools}{}
\makeatother


\title[Deklarativni operacijski sustavi]
{Deklarativni operacijski sustavi \\
Seminar}
\author[Karlo Pušelj]{Karlo Pušelj}
\institute[PMF--MO]{Prirodoslovno-matematički fakultet --- Matematički odsjek\\Sveučilište u Zagrebu}
\date{\today}

\begin{document}

\begin{frame}[plain]
	\titlepage
\end{frame}

\begin{frame}{Sadržaj}
	\tableofcontents
\end{frame}

\section{Uvod}

\begin{frame}{Motivacija}
	\pitanje{Što nakon kvara računala?}
\end{frame}

\begin{frame}{Imperativni i deklarativni pristup}
	\vskip25pt
	{\small
		\renewcommand{\arraystretch}{1.3}
		\begin{tabular}{>{\raggedright\arraybackslash}p{5.0cm} >{\raggedright\arraybackslash}p{5.0cm}}
			\toprule
			\textbf{Imperativni pristup}                                  & \textbf{Deklarativni pristup}             \\
			\midrule
			\enquote{Instaliraj SSH pa uključi servis.}                   & \enquote{SSH servis treba biti uključen.} \\
			Pamtimo što smo napravili.                                    & Zapisujemo što želimo imati.              \\
			Rezultat ovisi o početnom stanju i njegovim transformacijama. & Alat iz opisa priprema konfiguraciju.     \\
			\bottomrule
		\end{tabular}}
\end{frame}

\begin{frame}[fragile]{Kako to izgleda u NixOS-u?}
	\begin{lstlisting}[language=nix]
{
  services.openssh.enable = true;
}
\end{lstlisting}
	\begin{itemize}
		\item socijalni element: datoteku možemo čuvati u Git repozitoriju i dijeliti s drugima
	\end{itemize}
\end{frame}

\section{Kratki izlet u prošlost}

\begin{frame}{Kratki izlet u prošlost}
	\pitanje{Kako su tržišni i licencni odnosi oblikovali razvoj softvera?}

	\begin{itemize}
		\item konkurentski proizvođači UNIX-a razvijali su nekompatibilne varijante sustava \textrightarrow{} fragmentacija
		\item vlasničke licence ograničavaju dijeljenje i preraspodjelu softvera
		\item komercijalizacija sistemskog softvera narušava suradnju među programerima
	\end{itemize}

	\vfill
	{\scriptsize Izvor: \emph{The Open Group --- povijest UNIX-a} \cite{unix-history}; \emph{Richard Stallman --- GNU Manifest (1985.)} \cite{gnu-manifesto}}
\end{frame}

\section{Upravitelji paketa}

\begin{frame}{Upravitelji paketa}
	\pitanje{Što je upravitelj paketa?}
\end{frame}

\begin{frame}[fragile]{Što je upravitelj paketa?}
	Alat koji instalira, nadograđuje i uklanja softverske pakete te prati njihove ovisnosti.

	\vspace{0.5em}
	\begin{lstlisting}
aplikacija → biblioteka → druga biblioteka
\end{lstlisting}

	\begin{itemize}
		\item primjeri: \textbf{APT}, \textbf{DNF}, \textbf{pacman}
		\item paket sadrži datoteke programa i metapodatke
		\item ovisnosti opisuju što je programu potrebno za rad
	\end{itemize}
\end{frame}

\begin{frame}{Gdje tradicionalni model počne pucati?}
	\begin{itemize}
		\item višestruke verzije paketa ne mogu koegzistirati \textrightarrow{} \enquote{DLL hell}
		\item nadogradnje nisu atomičke
		\item nedostatak predvidljivosti
		\item nedostatak pratljivosti
		\item gradnja paketa je ovisna o stanju sustava \textrightarrow{} skrivene ovisnosti
		\item problemi s ažuriranjem korisnički modificiranih komponenti
	\end{itemize}
\end{frame}

\section{Nix, Nixpkgs i NixOS}

\begin{frame}{Nix, Nixpkgs i NixOS}
	\vskip25pt
	\centering
	{\renewcommand{\arraystretch}{1.3}
		\begin{tabular}{@{}l >{\raggedright\arraybackslash}p{8.0cm}@{}}
			\toprule
			\textbf{Naziv}   & \textbf{Što označava?}                                                  \\
			\midrule
			\textbf{Nix}     & Upravitelj paketa i sustav gradnje; također naziv jezika za opise.      \\
			\textbf{Nixpkgs} & Zbirka opisa paketa i biblioteka za rad s njima.                        \\
			\textbf{NixOS}   & Linux distribucija koja Nix koristi za izgradnju konfiguracije sustava. \\
			\bottomrule
		\end{tabular}}
\end{frame}

\begin{frame}[fragile]{Temelj Nix-a: Nix store}
	\begin{lstlisting}
/nix/store/<hash-A>-biblioteka-1.0/
/nix/store/<hash-B>-biblioteka-2.0/

program A → biblioteka-1.0
program B → biblioteka-2.0
\end{lstlisting}
	\begin{itemize}
		\item paketi se spremaju pod zasebne putanje u \texttt{/nix/store}
		\item postojeći rezultati gradnje ne mijenjaju se pri nadogradnji
		\item različite verzije i varijante mogu koegzistirati
	\end{itemize}
\end{frame}

\section{Derivacije i determinizam}

\begin{frame}[fragile]{Derivacije}
	\begin{lstlisting}
izvorni kod + ovisnosti + alat i upute za gradnju
                        ↓
                    derivacija
                        ↓
                 rezultat u storeu
\end{lstlisting}
	\begin{itemize}
		\item derivacija opisuje konkretan postupak gradnje i njegove ulaze
		\item ovisnosti zajedno tvore graf gradnje
		\item "binarna pričuvna memorija"
	\end{itemize}
\end{frame}

\begin{frame}{Postoji li ispravan redoslijed gradnje?}
	Modeliramo gradnju konačnim usmjerenim grafom $G = (V, E)$:
	\begin{itemize}
		\item vrh predstavlja korak gradnje
		\item brid $u \to v$ znači da gradnja $v$ treba rezultat gradnje $u$
	\end{itemize}

	\vspace{0.4em}
	\textbf{Teorem o topološkom poretku.} Konačan usmjeren graf ima topološki poredak ako i samo ako nema usmjerenih ciklusa.

	Topološki poredak je niz svih vrhova $v_1, \ldots, v_n$ takav da
	\[ v_i \to v_j \;\Longrightarrow\; i < j. \]

	\begin{itemize}
		\item svaku ovisnost možemo izgraditi prije paketa koji je koristi
		\item poredak ne mora biti jedinstven
	\end{itemize}

	\vfill
	{\scriptsize Izvor: \emph{Cornell --- topološki poredak} \cite{cornell-topological-sort}}
\end{frame}

\begin{frame}{Skica dokaza}
	Dokazujemo indukcijom po broju vrhova:
	\begin{itemize}
		\item[1.] prazan graf ima trivijalan poredak
		\item[2.] svaki neprazan konačan aciklički graf ima vrh bez ulaznih bridova; u suprotnom bismo neprestanim praćenjem prethodnika zbog konačnosti ponovili vrh i dobili ciklus
		\item[3.] odaberemo takav vrh $v$ i uklonimo ga; preostali graf i dalje je acikličan pa po indukcijskoj pretpostavci ima topološki poredak
		\item[4.] stavimo $v$ na početak tog poretka jer nema ulaznih bridova, a svi njegovi izlazni bridovi vode prema kasnijim vrhovima
	\end{itemize}

	\vspace{0.4em}
	Obratno, ciklus bi zahtijevao da se njegov početni vrh nalazi prije samoga sebe, što je nemoguće. \hfill $\square$
\end{frame}

\begin{frame}[fragile]{Malo o Nix jeziku}
	\begin{lstlisting}[language=nix]
a: a + 1 # lambda funkcija
a: b: a + b # lambda funkcija koja vraća lambda funkciju
{a = 1; b = 1} # atributni skup
{a, b}: a + b # lambda funkcija koja prima atributni skup
\end{lstlisting}
\end{frame}

\begin{frame}[fragile]{Malo o Nix jeziku (nastavak)}
	\begin{lstlisting}[language=nix]
{pkgs, nazivPosluzitelja, ...}: {
  networking.hostName = nazivPosluzitelja;
  services.openssh.enable = true;
  environment.systemPackages = [ ];
}
\end{lstlisting}
	\begin{itemize}
		\item \texttt{\{ \ldots \}}: skup atributa, odnosno parova ime--vrijednost
		\item \verb|"seminar"| je tekst, a \texttt{true} logička vrijednost
		\item \texttt{[ \ldots ]} je lista; elementi se odvajaju razmacima
		\item točke označavaju ugniježđene atribute; dodjelu završava \texttt{;}
	\end{itemize}
\end{frame}

\begin{frame}{Primjer derivacije iz Nixpkgs}
	\demo{pokazati \texttt{src} direktno iz Nixpkgs}
\end{frame}

\begin{frame}{Determinizam derivacija}
	\textbf{Što kontroliramo?}
	\begin{itemize}
		\item verziju paketa, izvore i ostale ulaze
		\item platformu za koju gradimo
		\item okruženje gradnje \textrightarrow{} pješčanik
	\end{itemize}
\end{frame}

\begin{frame}{Za raspravu}
	\pitanje{Kada derivacija ne bi bila deterministička?}
\end{frame}

\section{Praktični primjer}

\begin{frame}[fragile]{Praktični primjer: jednostavni server}
	\begin{lstlisting}[language=nix]
{ config, pkgs, ... }:

{
  # SSH pristup
  services.openssh.enable = true;

  # Vatrozid propušta samo SSH
  networking.firewall.allowedTCPPorts = [ 22 ];

  # Osnovni korisnik
  users.users.server.isNormalUser = true;
}
\end{lstlisting}
\end{frame}

\begin{frame}[fragile]{Generacije konfiguracija}
	\begin{lstlisting}
generacija 1 → generacija 2 → generacija 3
konfiguracija 1 → konfiguracija 2 → konfiguracija 3
\end{lstlisting}

	\demo{demonstracija generacija}

	\begin{itemize}
		\item nova se konfiguracija prvo izgradi
		\item prebacivanje profila na novu generaciju je atomično
		\item stara generacija ostaje dostupna dok je ne uklonimo
	\end{itemize}
\end{frame}

\begin{frame}[fragile]{Kako se izgradi sustav}
	\begin{lstlisting}
sudo nixos-rebuild switch
\end{lstlisting}

	\begin{center}
		\small opis \textrightarrow{} provjera i evaluacija \textrightarrow{} gradnja/preuzimanje \textrightarrow{} aktivacija
	\end{center}

	{\small
	\renewcommand{\arraystretch}{1.2}
	\begin{tabular}{@{}l >{\raggedright\arraybackslash}p{7.0cm}@{}}
		\toprule
		\textbf{Naredba \texttt{nixos-rebuild \ldots}} & \textbf{Rezultat}                                       \\
		\midrule
		\texttt{build}                                 & Izgradi konfiguraciju, bez aktivacije.                  \\
		\texttt{test}                                  & Aktivira je, bez postavljanja kao zadane za pokretanje. \\
		\texttt{boot}                                  & Postavi je za sljedeće pokretanje, bez aktivacije sada. \\
		\texttt{switch}                                & Aktivira je i postavi za sljedeće pokretanje.           \\
		\bottomrule
	\end{tabular}}
\end{frame}

\begin{frame}[fragile]{Što ako nova konfiguracija ne radi?}
	\begin{lstlisting}
sudo nixos-rebuild switch --rollback
\end{lstlisting}

	\begin{itemize}
		\item u pokrenutom sustavu možemo se vratiti na prethodnu generaciju
		\item ako se sustav ne pokreće, biramo sačuvanu generaciju u izborniku za pokretanje
		\item povrat konfiguracije ne vraća osobne datoteke, no njihov se povrat može riješiti pomoću BTRFS-a
	\end{itemize}

	\vfill
	{\scriptsize Izvor: \emph{povrat NixOS konfiguracije} \cite{nixos-manual-rollback}}
\end{frame}

\begin{frame}{Reproducibilnost}
	\pitanje{Hoće li konfiguracija generirati isti sustav danas i za mjesec dana?}
\end{frame}

\begin{frame}{Pahuljice (eng.\ \textit{flake}) spašavaju}
	\vskip25pt
	{\small
		\renewcommand{\arraystretch}{1.3}
		\begin{tabular}{@{}l >{\raggedright\arraybackslash}p{7.5cm}@{}}
			\toprule
			\textbf{Datoteka}          & \textbf{Uloga u projektu s flakes}           \\
			\midrule
			\texttt{configuration.nix} & NixOS postavke.                              \\
			\texttt{flake.nix}         & Deklarira ulaze i izlaze projekta.           \\
			\texttt{flake.lock}        & Bilježi zaključane revizije i hashove ulaza. \\
			\bottomrule
		\end{tabular}}

	\vspace{0.3em}
	\vskip10pt
	\begin{itemize}
		\item za ponavljanje konfiguracije dijelimo i zaključane ulaze
		\item nadogradnja ulaza je namjerna promjena koju možemo pregledati
	\end{itemize}
\end{frame}

\begin{frame}[fragile]{Korist i bez prelaska na NixOS}
	Primjer datoteke \texttt{shell.nix}:
	\begin{lstlisting}[language=nix]
{ pkgs ? import <nixpkgs> { } }:
pkgs.mkShell {
  packages = [ pkgs.git pkgs.curl ];
}
\end{lstlisting}

	\begin{lstlisting}
nix-shell
\end{lstlisting}

	\begin{itemize}
		\item dobivamo razvojnu ljusku s navedenim alatima
		\item izlaskom iz ljuske prestajemo koristiti to okruženje
		\item \texttt{<nixpkgs>} ovdje koristi lokalno odabranu zbirku; primjer \textbf{ne zaključava verzije}
	\end{itemize}
\end{frame}

\begin{frame}{Praktični primjeri}
	\demo{pokazati nekoliko praktičnih primjera}
\end{frame}

\begin{frame}{Zašto se Nix/NixOS ipak isplati?}
	\begin{itemize}
		\item zajednička konfiguracija za više računala
		\item sustav otvorenog koda s najširom podrškom i najvećim brojem paketa
		\item rad s cijelim operacijskim sustavom i paketima bez glavobolja
		\item čist sustav bez naknadnog smeća ako imate OKP
		\item praktičan rad s računalom
		\item lakša suradnja i dijeljenje
	\end{itemize}
\end{frame}

\begin{frame}{Pitanja i daljnje čitanje}
	\pitanje{Koji biste dio svojeg računala prvi opisali deklarativno?}

	\begin{itemize}
		\item {\footnotesize\url{https://nix.dev/}} --- vodiči i primjeri \cite{nixdev}
		\item {\footnotesize\url{https://nixos.org/manual/nixos/stable/}} --- NixOS priručnik \cite{nixos-manual}
		\item {\footnotesize\url{https://search.nixos.org/options}} --- pretraživanje NixOS opcija \cite{nixos-search-options}
		\item {\footnotesize\url{https://search.nixos.org/packages}} --- pretraživanje paketa \cite{nixos-search-packages}
	\end{itemize}
\end{frame}

\begin{frame}[allowframebreaks]{Literatura}
	\printbibliography[heading=none]
\end{frame}

\end{document}
